\documentclass{article}

\usepackage{amsmath,amssymb}
\usepackage{xurl}
\usepackage[hidelinks]{hyperref}
\setlength{\emergencystretch}{2em}

% Shared commands for notes in docs/paper-gaps/.

\DeclareUnicodeCharacter{2080}{\ifmmode{}_{0}\else\textsubscript{0}\fi}
\DeclareUnicodeCharacter{2081}{\ifmmode{}_{1}\else\textsubscript{1}\fi}
\DeclareUnicodeCharacter{2082}{\ifmmode{}_{2}\else\textsubscript{2}\fi}
\DeclareUnicodeCharacter{2099}{\ifmmode{}_{n}\else\textsubscript{n}\fi}

\newcommand{\Lean}{\textsc{Lean}}
\ExplSyntaxOn
\NewDocumentCommand{\leanid}{m}{
  \ifmmode
    \textsf{\detokenize{#1}}
  \else
    \group_begin:
    \urlstyle{sf}
    \tl_set:Nn \l_tmpa_tl {#1}
    \tl_replace_all:Nnn \l_tmpa_tl {\_} {_}
    \exp_args:NV \path \l_tmpa_tl
    \group_end:
  \fi
}
\ExplSyntaxOff
\providecommand{\tr}{\operatorname{tr}}
\newcommand{\tnleanIssueBase}{https://github.com/LionSR/TNLean/issues/}
\newcommand{\tnleanPrBase}{https://github.com/LionSR/TNLean/pull/}
\newcommand{\tnleanDocBase}{https://sirui-lu.com/TNLean/docs/}
\newcommand{\ghissue}[1]{\href{\tnleanIssueBase#1}{GitHub issue \##1}}
\newcommand{\ghpr}[1]{\href{\tnleanPrBase#1}{PR \##1}}
\newcommand{\doclink}[1]{\href{\tnleanDocBase#1.html}{\path{#1}}}

% Dirac notation and the identity operator, as in blueprint/src/macros/common.tex.
% \providecommand keeps note-local definitions authoritative.
\usepackage{dsfont}
\providecommand{\ket}[1]{|#1\rangle}
\providecommand{\bra}[1]{\langle #1|}
\providecommand{\braket}[2]{\langle #1 | #2 \rangle}
\providecommand{\ketbra}[2]{|#1\rangle\!\langle #2|}
\providecommand{\Id}{\mathds{1}}
\providecommand{\one}{\mathds{1}}
\providecommand{\C}{\mathbb{C}}
\providecommand{\MN}[1]{M_{#1}(\C)}
\providecommand{\MD}{M_D(\C)}

% External referencing for paper sources.  The build helper writes sanitized
% auxiliary files under build/paper-aux/ and excludes duplicated source labels.
\usepackage{xr}

% Import a generated paper auxiliary file with a caller-supplied label prefix.
% The candidate paths support builds from docs/paper-gaps/, the repository root,
% and build/paper-aux/ itself.
\newcommand{\importpaperaux}[2]{%
  \IfFileExists{../../build/paper-aux/#2.aux}{%
    \externaldocument[#1]{../../build/paper-aux/#2}%
  }{%
    \IfFileExists{build/paper-aux/#2.aux}{%
      \externaldocument[#1]{build/paper-aux/#2}%
    }{%
      \IfFileExists{#2.aux}{%
        \externaldocument[#1]{#2}%
      }{}%
    }%
  }%
}

% General external reference macros
\newcommand{\paperref}[2]{\ref{#1-#2}}
\newcommand{\papereqref}[2]{(\ref{#1-#2})}

% CPSV16 source (1606.00608 / MPDO-22-12-17-2.tex) external document and shortcuts
\importpaperaux{cpsv16-}{1606.00608/MPDO-22-12-17-2-tnlean}
\newcommand{\cpsvref}[1]{\paperref{cpsv16}{#1}}
\newcommand{\cpsveqref}[1]{\papereqref{cpsv16}{#1}}

% Verdict marker: \gapnote{<kind>}{<status>}. Kind is the policy
% classification (clarification, local-correction, scope-restriction,
% unfaithful, false-source, open-gap); status is open, wip (elimination
% underway), resolved, or historical. Machine-read by the site index;
% prints a verdict line.
\newcommand{\gapnote}[2]{\par\noindent\textbf{Verdict:} #1 (#2).\par}


\title{Jordan-Block Norm Estimate: the Shift Norm in One Dimension}
\author{}
\date{2026-08-13}

\begin{document}
\maketitle

\section{Source passage}

Fix an integer $D\ge 1$, let $N\in\MN{D}$ be the strict upper shift
$N_{i,j}=1$ if $j=i+1$ and $N_{i,j}=0$ otherwise, let
$J_{D}(\lambda)=\lambda\Id+N$ be the Jordan block with eigenvalue
$\lambda\in\C$, and write $\|\cdot\|_{\infty}$ for the largest singular value.
Chapter~8 of \cite{Wolf2012Quantum} derives, for every $n\in\mathbb N$ and
every natural number $k_{0}\le\min\{n,D-1\}$, the two-sided estimate
\begin{align}
  |\lambda|^{\,n-k_{0}}\binom{n}{k_{0}}
    \le \bigl\|J_{D}(\lambda)^{n}\bigr\|_{\infty}
    \le \sum_{k=0}^{\min\{n,D-1\}}|\lambda|^{\,n-k}\binom{n}{k}
\end{align}
of Equation~(8.104), at lines 1197--1215 of the local source file. The upper
bound is obtained from the binomial expansion
$J_{D}(\lambda)^{n}=\sum_{k}\binom{n}{k}\lambda^{n-k}N^{k}$ by the triangle
inequality, using the assertion $\|N\|_{\infty}=1$ for the strict upper shift
$N$ at line~1215.

\section{The deviation}

The equality $\|N\|_{\infty}=1$ holds for $D\ge 2$ but fails at $D=1$: the
strict upper shift on a one-dimensional space is the zero matrix, so
$\|N\|_{\infty}=0$. Wolf states the equality without a dimension restriction.
The Jordan block itself is defined for every $D\ge 1$, and Equation~(8.104) is
asserted for every such $D$, so the intermediate assertion is one degenerate
case wider than its proof supports.

\section{Formalized interpretation}

Only the direction $\|N\|_{\infty}\le 1$ enters the derivation: the triangle
inequality is followed by
$\|N^{k}\|_{\infty}\le\|N\|_{\infty}^{k}\le 1$, which is what discards the
shift factors from the expansion. The formalization therefore states
\leanid{Matrix.l2_opNorm_nilpotentShift_le_one} as the inequality
$\|N\|_{\infty}\le 1$, valid for every $D$, and derives
\leanid{Matrix.l2_opNorm_nilpotentShift_pow_le_one} from it. The proof is the
one Wolf indicates: the Gram identity
$N^{\dagger}N=\operatorname{diag}(0,1,\dots,1)$, which is
\leanid{Matrix.conjTranspose_nilpotentShift_mul_self}, gives a matrix whose
largest singular value is at most $1$, and the C*-identity
$\|N\|_{\infty}^{2}=\|N^{\dagger}N\|_{\infty}$ transfers the bound.

At $D=1$ the two sides of Equation~(8.104) coincide: the hypothesis
$k_{0}\le\min\{n,D-1\}$ forces $k_{0}=0$, and both bounds read
$|\lambda|^{n}\le\|\lambda^{n}\|_{\infty}\le|\lambda|^{n}$. The theorems
\leanid{Matrix.le_l2_opNorm_jordanBlock_pow},
\leanid{Matrix.l2_opNorm_jordanBlock_pow_le} and
\leanid{Matrix.l2_opNorm_jordanBlock_pow_bounds} therefore carry Wolf's
hypothesis set unchanged, with no dimension restriction beyond $D\ge 1$, which
the Jordan block presupposes.

\section{Verdict}

\textbf{Verdict: resolved local fix (shift norm at $D=1$).} Replacing the
source's intermediate equality $\|N\|_{\infty}=1$ by the inequality
$\|N\|_{\infty}\le 1$ removes a false degenerate case without weakening
Equation~(8.104), which is formalized with the source's own hypotheses.

\bibliographystyle{alpha}
\bibliography{references}

\end{document}
